\section{Threats To Validity}\label{sec:threats}
In this section, we report the possible threats to the validity of the results and conclusions presented in this study according to the classification framework proposed by Wohlin \etal \cite{wohlin_experimentation_2012}.
% \TODO{Report about each type of threat to the validity of the experiment, according to the classification framework proposed by Wohlin \etal \cite{wohlin12}.
% }

\subsection{Internal Validity}\label{sec:internal_threats}
% \TODO{Influences that can affect the dependent variable
% with respect to causality, without the researcher’s knowledge: history, maturation, testing, instrumentation, statistical regression, selection, mortality, ambiguity about direction of casual influence, interactions with selection, diffusion of imitation of treatments, compensatory equalization of treatments, compensatory rivalry, resentful demoralization.} 
Internal validity can be threatened through influence by external factors. For the ground truth collection, this threat was mitigated as much as possible through random treatment selection to mitigate order-related bias. The threat of carryover effects from one run affecting the next was mitigated by adding warm-up and cooldown periods between runs. For the simulator testing, order-related bias and carryover effects are only a threat to the measuring of the mean simulation time. These threats were mitigated by having a large sample size of simulation runs and simulating one repetition per treatment with all simulators before repeating with another repetition. 

Although Scaphandre is a widely used energy monitoring tool, it can cause data inconsistencies. Because of this, each treatment was run 30 times and the four repetitions with the energy consumption measures closest to the mean of these 30 repetitions were selected as traces to be simulated. As shown in \autoref{sec:accuracy}, some of the tested simulators are non-deterministic. In order to get an indication of the general performance of these simulators, every simulation of a workload was also run 30 times and the average over these 30 runs was taken.

\rone{We did not perform simulator-specific parameter calibration beyond configuring the hardware-dependent idle and maximum power consumption values required by the simulators. Introducing simulator-specific tuning would have raised additional threats to internal validity. In particular, the calibration process would have been performed by the authors, potentially introducing experimenter bias through subjective choices regarding parameter values and optimization strategies. Moreover, the simulators expose substantially different energy models and abstractions, resulting in large and heterogeneous calibration design spaces. Defining a representative and comparable calibration strategy across all simulators would therefore be challenging and could itself influence the observed results. For these reasons, we opted to evaluate each simulator using its default configuration and built-in power models, thereby ensuring a fair and reproducible comparison of their out-of-the-box capabilities.}

\rone{Evaluated simulators have heterogeneous support for communication modeling. While NetworkCloudSim and SimGrid explicitly model network communication and message exchanges, the other simulators provide limited or no communication-simulation capabilities. Since communication activities can affect resource utilization and, consequently, host-level energy consumption, differences in communication modeling may have influenced our results. Future work should explicitly assess the impact of communication simulation on energy-estimation accuracy.}

\subsection{External Validity}
% \TODO{Conditions that limit our ability to generalize the results of our experiment to, for example, industrial practice: interaction of selection and treatment, interaction of setting and treatment, interaction of history and treatment.}
External validity concerns the generalizability of the results discussed in this study to other systems and applications. Both \textit{Sock Shop} and \textit{Train Ticket} are microservice applications that are widely used in research \cite{fischer_overview_2024} (e.g., \cite{gias_atom_2019, wu_microrca_2020, ikram_root_2022, zhou_benchmarking_2018, zhang_deeptralog_2022}) and are of varying complexity. They are, however, academic benchmarking systems and may not be representative of other microservice systems that are more complex, have dynamic workloads or run on other technologies. Furthermore, the systems that were used to run the applications on were relatively small, which may not represent larger-scale microservice deployments. Finally, in this study we only look at the host energy usage of the cloud systems microservices run on. While this forms the biggest energy usage of microservices and other software systems, it does not cover the whole scope. Servers (including storage) are estimated to use 50\% of energy in a data centre. For networks, this is 8\% \cite{buyya_energy_2024}. As the other components such as cooling systems, UPS and PDU were not part of the simulations carried out for this study or the experiments to obtain the ground truth meaures, this means that approximately 42\% of energy usage of data centres was not taken into account. \textit{SimGrid} is the only simulator discussed here that has an option for modelling cooling systems with the Chiller plug-in, but this plug-in was not studied. Because of this, the accuracy results presented in this study can not be generalized to the suitability of these simulators for simulating overall energy usage of microservice applications, but only the energy usage of the application hosts.

\subsection{Construct Validity}
% \TODO{Generalizing the result of the experiment to the concept
% or theory behind the experiment: inadequate pre-operational explication of constructs, mono-operation bias, mono-method bias, confounding constructs and levels of constructs, interaction of different treatments, interaction of testing and treatment, restricted generalizability across constructs, hypothesis guessing, evaluation apprehension, experimenter expectancy.}
Construct validity is threatened by the executing of the experiment influencing the results. To mitigate this threat, the metric collection during the ground truth collection phase was done on a separate server. The monitoring tools such as Scaphandre and cAdvisor may have had an influence on the overall measured host power, but as all system properties such as idle power usage, maximum power usage and IPC were collected with these systems running as well, the influence should be minimized. Another threat is that for each metric only one tool was used, Scaphandre is widely used in research and peer-reviewed (e.g., \cite{centofanti_impact_2024, jay_experimental_2023, cadorel_protocol_2024, shankaranarayanan_poet_2024, gudepu_demonstrating_2024}), while cAdvisor is widely adopted in both research and industry. In order to mitigate the threat of the initial collection of cloud simulators not being broad enough and missing, for example, recently published simulators, a research practitioner from a relevant research background was consulted as a source for simulators as well. 

\subsection{Conclusion Validity}
% \TODO{Issues that affect the ability to draw the correct conclusion about relations between the treatment and the outcome of an experiment: low statistical power, violated assumption of statistical tests, fishing and the error rate, reliability of measures, reliability of treatment implementation, random irrelevancies in experimental setting, random heterogeneity of subjects.}
Conclusion validity concerns the reliability of the conclusions drawn in this study. To ensure this, all measurements were taken with 30 repetitions to increase the reliability of the results. Mean relative error was used as a metric for accuracy as it is widely used to evaluate the accuracy of scalar values while keeping the metric in the context of the domain. Whether there is evidence of significant differences between the simulators and power models was determined using the Mann-Whitney U and Kruskal-Wallis tests. For the second research question concerning the usability of the simulators, there is a threat to the validity of the results as these were only collected by one person. In order to mitigate this threat, the maintainers of the simulators were contacted and asked to confirm those results. However, only the maintainers for \textit{CloudSim}, \textit{NetworkCloudSim} and \textit{SimGrid} responded. Their feedback was taken into account when revising the paper.